\documentclass[UTF8,a4paper,11pt]{ctexart}
\usepackage{amsmath,amssymb}
\usepackage{geometry}
\usepackage{booktabs}
\usepackage{hyperref}
\geometry{margin=2.5cm}
\hypersetup{colorlinks=true,linkcolor=blue,urlcolor=blue}

\title{\textbf{比奈方程（Binet Equation）的推导\\及其在行星运动中的应用}}
\author{本地模型 Strata\\\small（Qwen3.8-Flash-Next，Gyro-S 量化）辅助生成}
\date{\today}

\begin{document}
\maketitle

\begin{abstract}
本文从极坐标下的牛顿第二定律出发，利用中心力场中的角动量守恒，
引入变量 \(u=1/r\)，逐步推导出比奈方程（Binet equation）。
随后将其应用于万有引力情形，证明行星轨道为圆锥曲线，
并讨论偏心率 \(e\) 与轨道类型的关系。
最后说明如何通过修改比奈方程的源项参数，
在数值模拟中改变行星的运动轨迹。
\end{abstract}

\tableofcontents
\newpage

\section{极坐标下的运动方程}

设质量为 \(m\) 的质点在中心力场中运动，其位置矢量用极坐标表示为
\begin{equation}
  \mathbf r = r\,\hat{\mathbf e}_r,
\end{equation}
其中 \(r\) 为径向距离，\(\theta\) 为极角。极坐标的单位矢量满足
\begin{equation}
  \dot{\hat{\mathbf e}}_r = \dot\theta\,\hat{\mathbf e}_\theta,
  \qquad
  \dot{\hat{\mathbf e}}_\theta = -\dot\theta\,\hat{\mathbf e}_r .
\end{equation}
对位置矢量求导，得到速度
\begin{equation}
  \mathbf v = \dot{\mathbf r}
  = \dot r\,\hat{\mathbf e}_r + r\dot\theta\,\hat{\mathbf e}_\theta .
\end{equation}
再次求导，得到加速度
\begin{equation}
  \mathbf a = \dot{\mathbf v}
  = \left(\ddot r - r\dot\theta^{2}\right)\hat{\mathbf e}_r
  + \left(2\dot r\dot\theta + r\ddot\theta\right)\hat{\mathbf e}_\theta .
\end{equation}
因此，牛顿第二定律 \(m\mathbf a=\mathbf F\) 在极坐标下的两个分量为
\begin{align}
  m\left(\ddot r - r\dot\theta^{2}\right) &= F_r, \label{eq:radial}\\
  m\left(2\dot r\dot\theta + r\ddot\theta\right) &= F_\theta . \label{eq:transverse}
\end{align}

\section{角动量守恒}

对中心力场，力只沿径向，横向分力为零：
\begin{equation}
  F_\theta = 0 .
\end{equation}
由式 \eqref{eq:transverse} 得
\begin{equation}
  m\left(2\dot r\dot\theta + r\ddot\theta\right)=0 .
\end{equation}
注意到
\begin{equation}
  \frac{d}{dt}\left(r^{2}\dot\theta\right)
  = 2r\dot r\dot\theta + r^{2}\ddot\theta
  = r\left(2\dot r\dot\theta + r\ddot\theta\right),
\end{equation}
于是
\begin{equation}
  \frac{d}{dt}\left(m r^{2}\dot\theta\right)=0 .
\end{equation}
定义角动量
\begin{equation}
  L \equiv m r^{2}\dot\theta = \text{常量}. \label{eq:L}
\end{equation}
由此可得角速度与半径的关系
\begin{equation}
  \dot\theta = \frac{L}{m r^{2}} . \label{eq:thetadot}
\end{equation}
这正是开普勒第二定律（面积定律）的微分形式。

\section{比奈方程的推导}

引入新变量
\begin{equation}
  u \equiv \frac{1}{r}, \qquad\text{即}\qquad r=\frac{1}{u}.
\end{equation}
记 \(u'=\dfrac{du}{d\theta}\)，\(u''=\dfrac{d^{2}u}{d\theta^{2}}\)。
由链式法则与式 \eqref{eq:thetadot}，
\begin{equation}
  \dot u = \frac{du}{d\theta}\dot\theta = u'\,\dot\theta ,
\end{equation}
而
\begin{equation}
  \dot r = -\frac{1}{u^{2}}\dot u
  = -\frac{u'}{u^{2}}\dot\theta .
\end{equation}
代入 \(\dot\theta = \dfrac{L}{m}u^{2}\)，得
\begin{equation}
  \dot r = -\frac{u'}{u^{2}}\cdot\frac{L}{m}u^{2}
  = -\frac{L}{m}u' . \label{eq:rdot}
\end{equation}
对式 \eqref{eq:rdot} 再求时间导数：
\begin{equation}
  \ddot r = -\frac{L}{m}\frac{d u'}{dt}
  = -\frac{L}{m}u''\dot\theta
  = -\frac{L}{m}u''\cdot\frac{L}{m}u^{2},
\end{equation}
即
\begin{equation}
  \boxed{\ \ddot r = -\frac{L^{2}}{m^{2}}u^{2}u''\ } \label{eq:rddot}
\end{equation}
同时
\begin{equation}
  r\dot\theta^{2}
  = \frac{1}{u}\left(\frac{L}{m}u^{2}\right)^{2}
  = \frac{L^{2}}{m^{2}}u^{3} . \label{eq:rthetadot2}
\end{equation}
将式 \eqref{eq:rddot} 与式 \eqref{eq:rthetadot2} 代入径向方程 \eqref{eq:radial}：
\begin{equation}
  m\left(-\frac{L^{2}}{m^{2}}u^{2}u''
  - \frac{L^{2}}{m^{2}}u^{3}\right)=F_r ,
\end{equation}
整理得
\begin{equation}
  -\frac{L^{2}}{m}u^{2}\left(u''+u\right)=F_r .
\end{equation}
两边同除以 \(-\dfrac{L^{2}}{m}u^{2}\)，得到\textbf{比奈方程}：
\begin{equation}
  \boxed{\
  \frac{d^{2}u}{d\theta^{2}} + u
  = -\frac{m\,F_r(1/u)}{L^{2}u^{2}}
  \ } \label{eq:binet}
\end{equation}
其中 \(F_r(1/u)\) 表示把径向力 \(F_r(r)\) 中的 \(r\) 用 \(1/u\) 替换。
式 \eqref{eq:binet} 把二维轨道问题化为一维关于 \(u(\theta)\) 的常微分方程。

\section{万有引力与圆锥曲线轨道}

对行星绕太阳运动，万有引力为
\begin{equation}
  F_r(r) = -\frac{GMm}{r^{2}}
  = -GMm\,u^{2},
\end{equation}
其中 \(G\) 为引力常量，\(M\) 为太阳质量，负号表示吸引力。
代入比奈方程 \eqref{eq:binet}：
\begin{equation}
  u''+u
  = -\frac{m\left(-GMm\,u^{2}\right)}{L^{2}u^{2}}
  = \frac{GMm^{2}}{L^{2}} .
\end{equation}
定义
\begin{equation}
  p \equiv \frac{L^{2}}{GMm^{2}} ,
  \label{eq:p}
\end{equation}
则方程化为常系数线性方程
\begin{equation}
  \boxed{\
  \frac{d^{2}u}{d\theta^{2}} + u = \frac{1}{p}
  \ } \label{eq:binet-newton}
\end{equation}
其通解为
\begin{equation}
  u(\theta) = \frac{1}{p}
  + A\cos\theta + B\sin\theta ,
\end{equation}
其中 \(A,B\) 为积分常数。令
\begin{equation}
  A = \frac{e}{p}\cos\theta_0,
  \qquad
  B = \frac{e}{p}\sin\theta_0 ,
\end{equation}
则
\begin{equation}
  u(\theta) = \frac{1}{p}\left[1+e\cos(\theta-\theta_0)\right],
\end{equation}
从而
\begin{equation}
  \boxed{\
  r(\theta) = \frac{p}{1+e\cos(\theta-\theta_0)}
  \ } \label{eq:conic}
\end{equation}
这正是圆锥曲线的极坐标方程，其中 \(p\) 为半通径，
\(e\ge 0\) 为偏心率，\(\theta_0\) 为轨道近心点方向。

\section{轨道类型与偏心率}

由式 \eqref{eq:conic} 可知，轨道类型完全由偏心率 \(e\) 决定：
\begin{center}
\begin{tabular}{ccl}
\toprule
偏心率 & 轨道类型 & 说明 \\
\midrule
\(e=0\) & 圆 & 半径 \(r=p\) 的圆周运动 \\
\(0<e<1\) & 椭圆 & 行星、多数卫星的轨道 \\
\(e=1\) & 抛物线 & 临界逃逸轨道 \\
\(e>1\) & 双曲线 & 非束缚的散射轨道 \\
\bottomrule
\end{tabular}
\end{center}
对于椭圆轨道，半长轴 \(a\) 与半通径 \(p\) 满足
\begin{equation}
  p = a\left(1-e^{2}\right),
\end{equation}
近日点与远日点距离分别为
\begin{equation}
  r_{\min}=\frac{p}{1+e}=a(1-e),
  \qquad
  r_{\max}=\frac{p}{1-e}=a(1+e).
\end{equation}
偏心率与总能量 \(E\) 的关系为
\begin{equation}
  e = \sqrt{1+\frac{2E L^{2}}{G^{2}M^{2}m^{3}}} .
\end{equation}
当 \(E<0\) 时 \(e<1\)，轨道为椭圆；当 \(E=0\) 时 \(e=1\)；
当 \(E>0\) 时 \(e>1\)。这体现了能量对轨道束缚性质的判据。

\section{推广与数值模拟}

比奈方程的价值在于：\textbf{给定任意中心力 \(F_r(r)\)，
都可以直接写出关于 \(u(\theta)\) 的轨道微分方程}，
而不必求解含时的二维运动方程。若将源项作如下推广：
\begin{equation}
  \frac{d^{2}u}{d\theta^{2}} + u
  = C + \lambda u^{2} + \mu u^{3},
  \qquad C=\frac{GMm^{2}}{L^{2}},
  \label{eq:general}
\end{equation}
则 \(\lambda,\mu\) 分别代表对平方反比定律的二次、三次修正。
当 \(\lambda=\mu=0\) 时回到牛顿引力，轨道为闭合圆锥曲线；
当修正项不为零时，轨道的近心点会随每圈转动，
形成\textbf{进动（precession）}或玫瑰线型轨道。
广义相对论对牛顿引力的最低阶修正恰好给出
一个正比于 \(u^{3}\) 的附加项，因此会自然导致水星近日点进动。

为直观展示上述结论，配套的交互网页 \texttt{index.html}
对式 \eqref{eq:general} 作四阶 Runge--Kutta 数值积分，
并实时绘制行星轨迹。用户可调参数包括
\begin{itemize}
  \item \(C\)：中心引力强度（对应 \(GMm^{2}/L^{2}\)）；
  \item \(e\)：初始偏心率（决定初始近心点距离）；
  \item \(\lambda\)：二次修正系数；
  \item \(\mu\)：三次修正系数（可模拟相对论进动）；
  \item 动画速度与轨迹显示。
\end{itemize}

\section*{参考文献}
\begin{enumerate}
  \item H. Goldstein, C. P. Poole, J. L. Safko,
        \emph{Classical Mechanics}, 3rd ed., Pearson, 2002.
  \item L. D. Landau, E. M. Lifshitz,
        \emph{Mechanics}, 3rd ed., Butterworth--Heinemann, 1976.
  \item S. Weinberg,
        \emph{Gravitation and Cosmology}, Wiley, 1972.
\end{enumerate}

\end{document}
